\section{Related Work}

\subsection{JEPA world models and latent planning}

LeCun's proposal for autonomous machine intelligence \cite{lecun-2022-path} places a
joint-embedding predictive architecture at the centre of a world model: predict in
representation space, not in observation space, and use the learned latent dynamics for
planning. DINO-WM \cite{zhou-2025-dinowm} showed that this works zero-shot on pre-trained
visual features, planning with either gradient descent or the cross-entropy method. Sobal et
al.\ \cite{sobal-2025-offline} argue that reward-free offline trajectories are sufficient for
planning with latent dynamics models. Wang et al.\ \cite{wang-2026-temporal} add a
temporal-straightening criterion that makes latent trajectories nearly linear, which
conditions gradient-based planning better; AdaJEPA \cite{wang-2026-adajepa} is built on top
of that recipe. LeWorldModel \cite{maes-2026-lewm} removes the remaining fragility by training
a JEPA end-to-end from pixels with only two loss terms---a next-embedding prediction loss and
a regulariser that pushes latent embeddings towards an isotropic Gaussian---with no
exponential moving average, no pre-trained encoder and no auxiliary supervision; it also uses
prediction surprise to detect physically implausible events. V-JEPA~2
\cite{assran-2025-vjepa2} demonstrates the same paradigm at foundation-model scale.

We adopt LeWorldModel's training philosophy, because it makes a from-scratch, single-machine,
end-to-end study feasible, and we use its surprise notion as the resource signal of our
adaptation controller.

\subsection{Test-time adaptation}

Test-time training and adaptation (TTT/TTA) update a pre-trained model on the test input using
an unsupervised objective: auxiliary rotation prediction \cite{sun-2020-ttt}, entropy
minimisation \cite{wang-2021-tent}, augmentation consistency \cite{zhang-2022-memo}, masked
reconstruction \cite{gandelsman-2022-mttt}, and stabilised variants for wild, non-stationary
streams \cite{niu-2022-efficient,niu-2023-stable,wang-2022-cotta}. The literature is almost
entirely about classifiers, where the adaptation signal is the model's own confidence. In
latent planning the situation is different: the environment hands the agent a labelled
transition $(o_t,a_t,o_{t+1})$ for free, so the adaptation objective is exactly the
pretraining objective, and it can be optimised in a closed loop with control.

AdaJEPA \cite{wang-2026-adajepa} is the closest work to ours and our main baseline: it adapts
inside the MPC loop, one step per replan, restricted to the predictor's last transformer block
and the encoder head, with a bounded replay buffer of the most recent transitions. In their
ablations the adaptation target is chosen by hand per environment family, and stronger
adaptation (larger learning rate, more steps, larger buffer) is sometimes beneficial
out of distribution---but only when latency permits, and at the cost of stability in
distribution. Our work replaces those hand-set knobs with an uncertainty-driven controller and
makes adaptation stateful across episodes.

\subsection{Adaptation in planning and control}

Online model adaptation is classical in adaptive control, from self-adapting IDCOM
\cite{foigel-1979-idcom} to online system identification in Dyna-style model-based RL
\cite{sutton-1991-dyna,thrun-1990-planning}. Modern world-model work adapts predictors with
latent action models \cite{gao-2025-adaworld}, residual latent-state dynamics
\cite{lanier-2025-residual}, or data synthesis during training to close the train--test gap
\cite{parthasarathy-2025-closing}; all of them require either target-domain fine-tuning data,
extra online rollouts, or an outer self-improvement loop. TD-MPC2 \cite{hansen-2024-tdmpc2}
learns a latent dynamics model and a terminal value function online but optimises a reward
objective rather than a self-supervised latent prediction. We stay strictly in the
reward-free, self-supervised regime.

\subsection{Recursive and hierarchical latent computation}

The Hierarchical Reasoning Model \cite{hrm-2025} couples a slow, high-level recurrent module
with a fast, low-level one, and lets the fast module iterate many times per slow update;
HRM-Text \cite{hrm-text-2026} scales that architecture to a 1B language model and shows that
the hierarchical recurrence substitutes for a large amount of pretraining compute. Tiny
Recursive Models \cite{tiny-recursive-2025} simplify the idea to its core: a
$\sim$7M-parameter network recursively improves a candidate answer in latent space, updating a
latent state and then the answer, with supervision applied at every improvement step, reaching
strong puzzle-solving performance without scale.

We take the two structural lessons---\emph{two coupled timescales} and \emph{recursion with
deep supervision}---and place them inside a world model's predictor. To our knowledge this is
the first recursive-refinement JEPA predictor, and the first use of refinement depth as an
adaptation-time compute knob.

\subsection{Continual learning and safety of online updates}

Catastrophic forgetting \cite{mccloskey-1989-forgetting} is usually mitigated with elastic
weight consolidation \cite{kirkpatrick-2017-ewc} or replay \cite{riemer-2019-replay}; in
streaming TTA, parameter restoration and sharpness-aware variants are used
\cite{niu-2023-stable,wang-2022-cotta}. Anchoring to a pretrained reference is the standard
anti-drift device. Our safety layer is deliberately minimal and auditable: a Euclidean anchor
ball on the fast parameters, a validation gate on a held-out slice of the replay buffer with
rollback, an immutable slow anchor, and a support penalty expressed as a Mahalanobis radius
computed from the training latents.

\subsection{Reproducibility-focused benchmarks}

PushT and PushObj \cite{chi-2025-diffusion-policy,zhou-2025-dinowm} and PointMaze
\cite{fu-2021-d4rl} are the standard test beds for latent planning, and the diverse-maze
protocol \cite{zhang-2026-hierarchical} standardises layout shift. Reproducing these requires
the original assets (expert datasets, pre-trained checkpoints and MuJoCo models); where those
are unavailable, a physics-based re-implementation that preserves the protocol is a useful
alternative. We build ours with numpy only---impulse-based contact physics and a software
rasteriser---so that the entire study, from data generation to planning evaluation, runs on a
laptop CPU with no external simulator, and so that every shift condition is generated by an
explicit, inspectable parameter.
